\documentclass[11pt]{article}
\usepackage[margin=1in]{geometry}
\usepackage{hyperref}
\usepackage{enumitem}
\usepackage{titlesec}
\usepackage{parskip}

\hypersetup{colorlinks=true, linkcolor=blue, urlcolor=blue}
\titlespacing*{\section}{0pt}{8pt}{4pt}
\titlespacing*{\subsection}{0pt}{6pt}{2pt}

\title{\vspace{-1.5cm}Bulk Mailing Tool \\ \large A Mail-Scheduling System for GYWS, IIT Kharagpur}
\author{Megha Singh}
\date{August 2025}

\begin{document}
\maketitle
\vspace{-0.8cm}

\section*{Objective}
GYWS is a student-run NGO at IIT Kharagpur that runs donor outreach and awareness campaigns over email.
Sending these campaigns manually from a single inbox is slow and trips spam filters, since providers flag
bursts of near-identical messages from one sender. The Bulk Mailing Tool was built to let a small team
schedule and send large recipient lists in a controlled, trackable way, reducing manual effort and
spam-filter risk.

\section*{What was built}
The system is split into two services:

\begin{itemize}[nosep]
  \item \textbf{Frontend} — a React (Vite) dashboard, styled with Tailwind CSS, for composing campaigns
  with a rich-text editor, uploading CSV recipient lists, managing sender ``agents,'' and reviewing send
  history.
  \item \textbf{Backend} — a FastAPI service backed by MongoDB. It exposes JWT-authenticated endpoints to
  create sender agents, send an email immediately, or schedule a single email or a full batch for future
  delivery. Scheduled jobs are handled by APScheduler and dispatched through \texttt{smtplib} over SMTP.
  Each recipient's message supports placeholder substitution (e.g.\ \texttt{\{\{name\}\}}) so a single
  template can be personalized per recipient, and every send is logged with its delivery status.
\end{itemize}

\section*{Scope note}
This write-up accompanies a CV entry that also references a Redis/Bull-based queue for per-user rate
limiting and a Mistral-AI-generated content-variation step to reduce duplicate-content spam flags. Those
pieces were part of the intended design and were evaluated during development; the codebase in this
repository reflects the FastAPI + APScheduler implementation that was actually shipped and used to send
GYWS's campaigns. The rate-limiting and content-variation ideas remain a natural next iteration on top of
the current scheduler.

\section*{Outcome}
The tool was used to send donor and outreach campaigns for GYWS, supporting multiple scheduled batches
with per-recipient personalization and delivery tracking, in support of the organization's fundraising
and awareness efforts during 2025.

\section*{Repositories}
\begin{itemize}[nosep]
  \item Frontend: \url{https://github.com/GYWS-TechOps/mailkaro-frontend}
  \item Backend: \url{https://github.com/GYWS-TechOps/mailkaro-backend}
  \item Consolidated repository: \url{https://github.com/MeghaSinghall/bulk-mailing-tool}
\end{itemize}

\end{document}
